\subsection*{Математическое ожидание}
Для начала, хотелось бы понять "--- а зачем это всё надо? Видите ли, разные выборки имеют разную природу, и заранее точно мы не можем знать, какую модель будет \textit{корректно} применить для получения хороших результатов. Соответственно, предлагается на основе ошибки, которая допускает модель, вывести дополнительные предположения о том, насколько хорошо справляется наша модель.

Например, рассмотрим случай регрессии с \texttt{MSE}. Также для простоты представим, что наша модель учится предсказывать зависимость вида
\[
y = f(x) + \epsilon
\]
Где $\epsilon$ "--- некоторый шум, искуственно добавленный в данные. В реальной жизни этим шумом может быть
\begin{itemize}
    \item Недостаточная точность измерительного прибора;
    \item Плохой набор данных, на основе которого невозможно в точности предсказать целевую переменную;
    \item Погрешность;
    \item ...
\end{itemize}

Как мы уже и обсуждали на одном из первых занятий, на основе одного лишь
\[
(a(x) - y)^2
\]
нельзя сделать выводы о том, насколько хороша построенная модель. Но можем ли мы как"=то сделать так, чтобы какие"=то выводы всё таки было возможно сделать?

Ну например, хотелось бы в модели смотреть не на конкретный объект, а на всю выборку, откуда этот объект поступил:
\[
a(x) = a(x, X)
\]
Пока представим, что мы как"=то переоборудовали нашу модель так, чтобы она на вход ещё принимала всю выборку.

\newpage

И здесь мы подходим к самому интересному: хотелось бы с такой моделью, которую мы построили ранее, увидеть \texttt{MSE} на \textit{всех возможных данных} (не только на тех, которые даны в $X$, а вообще на всевозможных таких данных). Идея в том, что с точки зрения статистики, $X$ является лишь частью так называемой \textit{генеральной выборки}. Но как нам получить доступ ко всей, \textit{генеральной} выборке?

Тут мы плавно переходим к понятию \textit{математического ожидания}. Давайте сразу взглянем на формулу:
\[
MSE = \mathbb{E}[(y - a(x,X))^2]
\]
Математическое ожидание обозначается буквой $\mathbb{E}$, и в контексте этой формулы означает: <<Если мы будем повторять этот процесс (строить модель, делать предсказание, считать ошибку) много раз на разных данных, какое среднее значение ошибки мы получим?>>

Математическое ожидание "--- это инструмент, который позволяет нам перейти от конкретных ошибок к общей картине. Вместо того чтобы говорить <<в этот раз ошибка 0.2>>, мы говорим <<в среднем, учитывая все возможные варианты обучающих данных, ошибка составляет...>>.